% \maketitle


\begin{abstract}
Stuxnet is one of the most influential computer viruses of the 2010’s and has played a big impact on the international cybersecurity industry. This is because it is known as the first nation state led cyberattack on another nation, which could have major political ramifications.[16] Stuxnet is believed to be developed over many years between the United States and Israel intelligence services with Iran's nuclear program as its target.[13] This program had the codename of Operation Olympic games and was originally started in 2006 during the Bush administration and continued during the Obama administration. This sophisticated computer virus would eventually escape its target of Iran's uranium enrichment efforts [18] Becoming a big talking point in the international IT/cybersecurity community as it started to affect computers outside Iran. Stuxnet would also have an impact on computing ethics due to it being the first nation state cyberattack on another nation. Leading it to be a turning point in computer ethics. While Stuxnet did do physical damage that damage did not include human injuries or death. Which means the ethical concerns are mainly on why and if it should have been used. One of the biggest factors of this is the reasoning behind Stuxnet was to take out Iranian nuclear capabilities through a cyberattack. Specifically to prevent Israel by taking conventional measures by using missiles to attack. Which would have led to many innocent deaths. Because of this the ethical considerations are a bit more profound and lean towards why this is more acceptable as the alternative would have surely led to injury and death. Previously there have been many articles on this topic over the past 15 years with a documentary also existing.
\end{abstract}


\section{Introduction}

Stuxnet is one of the most influential computer viruses of the 2010’s and has played a big impact on the international cybersecurity industry. This is because it is known as the first nation state led cyberattack on another nation, which could have major political ramifications.[16] Stuxnet is believed to be developed over many years between the United States and Israel intelligence services with Iran's nuclear program as its target.[13] This program had the codename of Operation Olympic games and was originally started in 2006 during the Bush administration and continued during the Obama administration. This sophisticated computer virus would eventually escape its target of Iran's uranium enrichment efforts [18] Becoming a big talking point in the international IT/cybersecurity community as it started to affect computers outside Iran. Stuxnet would also have an impact on computing ethics due to it being the first nation state cyberattack on another nation. Leading it to be a turning point in computer ethics. While Stuxnet did do physical damage that damage did not include human injuries or death. Which means the ethical concerns are mainly on why and if it should have been used. One of the biggest factors of this is the reasoning behind Stuxnet was to take out Iranian nuclear capabilities through a cyberattack. Specifically to prevent Israel by taking conventional measures by using missiles to attack. Which would have led to many innocent deaths. Because of this the ethical considerations are a bit more profound and lean towards why this is more acceptable as the alternative would have surely led to injury and death. Previously there have been many articles on this topic over the past 15 years with a documentary also existing.


\section{Case Summary}
Stuxnet was a virus discovered by a Belarusian antivirus company VirusBlokAda employee Sergey Ulasen whose customers In Iran the virus was attacking.[18] Stuxnet refers to a piece of malicious computer code which was discovered on June 17, 2010. Stuxnet specifically targeted supervisory control and data acquisition (SCADA) systems which are responsible for the management and control of hardware and industrial systems and is believed to be responsible for causing damage to the Iran nuclear program after its first installation at the Natanz Nuclear Facility in 2009. Once it was found It was given the name Stuxnet by cybersecurity specialists but was named "the bug" [13]  internally while it was being developed. The code of the virus was found to be 20 times the size of normal virus code. Due to this size and sophistication the consensus was that the worm was developed by a nation state and not a hacker group.(Kushner) Stuxnet was also the first malicious software (malware) designed specifically to attack a particular type of industrial control system. In this case the suspected target was industrial processes controlled by Siemens SIMATIC WinCC and PCS 7 control systems running on windows.[18] Stuxnet used 4 zero-day attacks to infect systems then it would utilize a kernel-mode rootkit to remain undetected on the infected windows system, this rootkit consisted of device drivers which were digitally signed with stolen certificates from Realtek and JMicron which allowed Stuxnet to appear as trusted software from well-known vendors allowing it to remain undetected for a long period of time. The code itself was known as a self-replicating worm that attacked programmable logic controllers in the uranium enrichment centrifuges. Stuxnet was first spread via USB and then used its capabilities as a worm to infect other devices. By targeting these centrifuges the virus could speed up or slow down how fast they were spinning and both would cause the same amount of damage to destroy the centrifuge. It would monitor the frequency of the motor attached and only attack systems that spin between 807 Hz and 1210 Hz and periodically modify the frequency to 1410 Hz then 2 Hz then 1064 Hz. Stuxnet reportedly destroyed over ⅕ of Iran’ nuclear centrifuges and infected over 200000 computers and caused 1000 machines to physically degrade. It also spread to many different countries with a study by Symantec showing that the main affected countries were Iran, Indonesia, and India with 58.9% of infected devices being Iranian, 18.2% being Indonesian, and 8.3% being Indian as well as small percentages of other countries like Azerbaijan, the U.S., and Pakistan. This attack would lead to many ethical and political conversations due to the nature of the attack and how it has affected relations between countries and cyberattacks. A future consideration is that “the example of Operation Olympics Games shows that complex cyber-sabotage operations resulting in the destruction of critical infrastructure on a large scale require the involvement of numerous state resources and advanced cyber activities, and the use of many different methods and intelligence sources. Thus, strong states with well-developed intelligence capabilities are much more capable of effectively using cyber-sabotage on a large scale.” [13]

\section{Ethical Analysis}
Act Utilitarianism:

Act Utilitarianism judges an action by the resulting consequences. An action is ethical if it results in more good over all than any other action which was available [1], [2]. Evaluating Stuxnet with this ethical theory involves comparing it with realistic alternatives which the United States and Israel could have chosen instead. The first alternative is a military airstrike on Iran’s nuclear facilities, which would result in large scale damage to the program but also risked the threat of a true armed conflict. The second alternative is relying entirely on diplomatic action and economic sanctions, which would avoid sabotage and force but would also leave the Iranian nuclear program fully intact. The tables below score how each option affects each stakeholder group, on a scale from -10 (severe harm) to +10 (major benefit). The resulting totals will show which option produces the greatest total good.

Stuxnet Malware Stakeholder Table:

\begin{table}[h]
\centering
\begin{tabular}{|c|c|p{8cm}|}
\hline
Stakeholder & Stuxnet & Reasoning \\ \hline
Iranian government/leadership & -6 & Iran lost around 1,000 centrifuges and an estimated 1-2 years of progress. The government also suffered international embarrassment but neither the government nor the program was destroyed. \\ \hline
Iranian nuclear scientists \& engineers & -3 & The work spent was wasted and they faced suspicion from the government since Iran initially blamed internal errors or sabotage. No physical harm was caused. \\ \hline
Iranian Civilians & -1 & There were almost no direct effects on daily life. \\ \hline
United States & +6 & Delayed Iranian uranium enrichment without a war or US casualties. The only downside was the loss of control once the code escaped. \\ \hline
Israel & +7 & Reduced the thread of Iran on Israel while avoiding any diplomatic problems or risk of retaliation. \\ \hline
Regional neighbors of Iran & +3 & The delay of Iran’s nuclear program reduced the pressure for a regional arms race while avoiding war. \\ \hline
Global bystanders & -3 & The worm spread far beyond Iran to thousands of machines worldwide causing cleanup costs and security concerns. \\ \hline
Global Public & -2 & Normalized state-sponsored cyberattacks against physical infrastructure and gave other malicious actors a framework to study. \\ \hline
Future victims of cyberattacks & -4 & Iran responded to Stuxnet by creating its own cyber taskforce. The Saudi Aramco attack and the US bank attacks in 2012 are widely seen as connected. \\ \hline
TOTALS: & -3 & Sets a precedent for cyberattacks but has no physical consequences. \\ \hline
\end{tabular}
\end{table}

The end result is negative but the negative consequences do not outweigh the positive benefits by much meaning that Stuxnet was a reasonably balanced action to take although still negative in its consequences.

Airstrike/Military Action Stakeholder Table:

\begin{table}[h]
\centering
\begin{tabular}{|c|c|p{8cm}|}
\hline
Stakeholder & Airstrike & Reasoning \\ \hline
Iranian government/leadership & -9 & Physical destruction of facilities would be a major military humiliation. However it would likely lead to support from the citizens. \\ \hline
Iranian nuclear scientists \& engineers & -8 & Strikes on facilities would likely kill or injure workers. \\ \hline
Iranian Civilians & -7 & Risk of casualties near sites and possible radiation contamination. Also a high risk of being involved in a war. \\ \hline
United States & +2 & Could cause a larger setback for the nuclear program than Stuxnet but at the cost of being involved in conflicts and attacks. \\ \hline
Israel & +4 & A large short-term setback for Iran but would cost international criticism. \\ \hline
Regional neighbors of Iran & -5 & Risk of being caught in a war with attacks on infrastructure and disruption of the Strait of Hormuz. \\ \hline
Global bystanders & 0 & No direct effect on computers or systems outside of the region. \\ \hline
Global Public & -4 & Preventative strikes on a foreign nation weaken norms against force and disruption of oil shipments would affect the economy. \\ \hline
Future victims of cyberattacks & 0 & No cyberattack precedent was created. \\ \hline
TOTALS: & -27 & Has severe political and ethical consequences. \\ \hline
\end{tabular}
\end{table}

The Airstrikes/Military Action alternative has severe negative consequences for almost all parties involved and even the benefits granted are very small meaning the consequences far outweigh the benefits.

Diplomacy and Sanctions Stakeholder Table:

\begin{table}[h]
\centering
\begin{tabular}{|c|c|p{8cm}|}
\hline
Stakeholder & Sanctions & Reasoning \\ \hline
Iranian government/leadership & -2 & Economic and diplomatic pressure but the nuclear program continues. \\ \hline
Iranian nuclear scientists \& engineers & 0 & Work would continue uninterrupted with some difficulty getting materials due to sanctions. \\ \hline
Iranian Civilians & -4 & Sanctions causing a rise in inflation and shortages of medicine and goods. \\ \hline
United States & -2 & No war but Iran gets to keep developing their nuclear program. \\ \hline
Israel & -5 & Iran’s program continues uninterrupted raising the perceived threat to Israel. \\ \hline
Regional neighbors of Iran & -2 & A more advanced nuclear program pushes neighbors to develop their own programs to keep up. \\ \hline
Global bystanders & 0 & No direct effect. \\ \hline
Global Public & +1 & Reinforces diplomacy and international institutions over force. \\ \hline
Future victims of cyberattacks & 0 & Sets no cyber precedent. \\ \hline
TOTALS: & -14 & Has no physical consequences but allows Iran’s nuclear program to develop as it would just with some restrictions. \\ \hline
\end{tabular}
\end{table}

Diplomacy and Sanctions have a negative score for almost all parties involved although not severe. There are very few benefits to putting sanctions on Iran as their nuclear program still grows and now the U.S. and Israel have to be directly involved.

By Act Utilitarianism [1], [2], Stuxnet had the least consequential result of the 3 actions which could have been chosen and as such was the most ethical choice to make. All of the options have their consequences but the worst of all was the U.S. or Israel taking military action such as an airstrike because in that scenario almost everyone involved would be negatively affected with immense consequences for Iran. The next best option would be diplomacy and sanctions which would mean no conflict takes place but there are far less benefits for either side of the deal as the U.S. or Israel would have to become directly politically involved in Iran’s affairs and Iran would have to deal with sanctions which would also negatively impact the citizens.

Rule Utilitarianism:
Rule Utilitarianism follows the idea that an action is ethical if it follows a rule which, if generally adopted, would produce the greatest overall good [2], [3]. Unlike act utilitarianism, it does not weigh the consequences of each individual action [3]. Instead, it asks which rules, if followed, lead to the best outcomes for society as a whole. This means that an action can be justified even if it causes some harm as long as the rule behind it leads to more good than harm when followed. To evaluate Stuxnet we will follow the rule “Nations should use targeted, non-lethal sabotage to prevent new powers from acquiring nuclear weapons, because every new nation with access to a nuclear arsenal increases the risk of catastrophic use.” The reasoning behind this rule is that every nuclear-armed state adds new chances for disaster. Even if each new state is unlikely to cause a catastrophe on its own, the risk accumulates with every added state. Because the consequences of nuclear war are so disastrous even a small reduction in that risk produces a large amount of good. This rule has two important limitations. First, the sabotage must be targeted, meaning that it is aimed at weapons programs rather than a population or general infrastructure. Second, it must be non-lethal which rules out bombing and assassinations. These limits keep the rule from justifying violence.

Under this rule, Stuxnet is strongly supported. It was designed specifically to target centrifuges at Iran’s Natanz enrichment facility, killed no one, and delayed Iran’s nuclear program by an estimated one to two years [15]. Stuxnet meets all of the conditions set by the general rule. Rule Utilitarianism also helps to justify the main criticism of Stuxnet which was the harm to bystanders caused by the unintended spread beyond Iran to thousands of computers around the world [15]. Rule Utilitarianism does not ask whether every action was good, it asks whether following the rule produces more good than harm overall [3]. Some collateral damage is an accepted cost of many rules such as self defense, for example sometimes causes unintended harm. By this logic the temporary consequences of Stuxnet spreading were small compared to the benefit of reducing the risk of nuclear war.

Virtue Ethics:
Virtue Ethics, rooted in Aristotle’s Nicomachean Ethics [4], judges an action by whether it reflects good character rather than by its consequences or fixed rules. Aristotle defines each virtue as a mean between two extremes, one of excess and one of deficiency [4, Bk. II]. Courage, for example, is the mean between cowardice which is too little willingness to face danger and recklessness which is too much [4, Bk. III]. The mean is not simply a midpoint, it is the right response for the specific situation and finding it requires the ability to see what the situation calls for and act on it [4, Bk. VI]. The United States and Israel had a range of potions available in response to Iran’s nuclear program. One extreme was inaction, relying on diplomacy and sanctions while Iran continued to enrich uranium and get closer to nuclear weapons. This can be seen as the extreme of deficiency or cowardice since it avoids all risk while leaving the threat unaddressed. At the other end of the extreme is a military airstrike on Iran’s uranium enrichment facilities. Although this action would have ceased Iran’s ability to develop nuclear weapons, it is at the risk of civilian casualties and retaliation. This can be seen as the reckless extreme since it accepts danger which is above what was needed. Stuxnet falls between these extremes, it took direct action to stop the threat by physically damaging the centrifuges and delaying the enrichment of uranium all while killing no one and avoiding war [15]. In this way it reflects courage as a virtue as well as other virtues such as temperance or restraint all while showing wisdom in finding an effective option without the cost of lives or inaction. From this perspective Stuxnet was the virtuous mean, a response matched to the situation.

\section{Recommendations for Resolving the Case}
Stance on the Historical Resolution

Looking back at how Stuxnet unfolded, the operation achieved its immediate tactical goal: it physically wrecked around 1,000 centrifuges at Natanz and set back Iran's enrichment efforts by one to two years without causing direct human casualties [13], [15], [16]. Compared to an outright military airstrike—which would have killed workers on-site and risked an armed regional conflict—Stuxnet looked like a restrained, surgical alternative [13]. However, treating the case as a success overlooks its long-term fallout.

From an Act Utilitarian perspective, the math falls apart once the code escaped the facility. Stuxnet infected over 100,000 machines across countries like Indonesia, India, and the U.S., burning four unpatched Windows zero-days and exposing unrelated networks to collateral risk [1], [2], [10], [15]. On top of that, Virtue Ethics and engineering codes like SECEPP focus on character and professional honesty [5], [12]. The weapon relied on active deception by spoofing normal sensor readouts to human operators while secretly destroying equipment [10]. Releasing an autonomous worm with flawed containment showed a deficit in the virtues of prudence and temperance, directly conflicting with professional obligations to protect public safety and disclose known hazards (SECEPP Clauses 1.03 and 1.04) [12]. Once the worm was discovered and decompiled, it essentially published a blueprint for attacking critical civilian utilities worldwide [15], [19].

What Would Have Been a More Ethical Resolution?

A more ethical resolution would have addressed the nuclear threat without compromising shared digital infrastructure and public trust:

- Strict Non-Propagating Cyber Action (If Software Was Used): If a cyber operation was chosen over kinetic strikes, the payload should never have included autonomous network replication modules like print spooler and SMB exploits [10]. A strictly targeted, closed-loop diagnostic or direct logic payload would have met Rule Utilitarian constraints by eliminating collateral spillover to thousands of civilian machines [3], [10], [15].
- Multilateral Diplomacy and Enhanced IAEA Inspections: A rules-based, deontological resolution would rely on strengthened International Atomic Energy Agency (IAEA) monitoring backed by targeted international economic sanctions [11]. While slower and requiring greater political patience, diplomacy avoids setting a precedent for state-sponsored cyber sabotage against industrial machinery, preserving international norms and stability without releasing dangerous exploits into the public domain [15], [19].
- Responsible Vulnerability Disclosure: Hoarding and weaponizing four zero-day vulnerabilities left commercial enterprises and municipal utilities exposed to opportunistic exploitation [10], [15]. A resolution aligned with computing ethics would require that once operational objectives were achieved, critical operating system vulnerabilities be responsibly disclosed to vendors, upholding SECEPP Clause 1.04 to minimize broader public harm [12].

\section{Conclusion}

%This is a test. \cite{ref01_bentham1907}

Looking back at how Stuxnet unfolded, the operation achieved its immediate tactical goal: it physically wrecked around 1,000 centrifuges at Natanz and set back Iran's enrichment efforts by one to two years without causing direct human casualties [13], [15], [16]. Compared to an outright military airstrike—which would have killed workers on-site and risked an armed regional conflict—Stuxnet looked like a restrained, surgical alternative [13]. However, treating the case as a success overlooks its long-term fallout.

From an Act Utilitarian perspective, the math falls apart once the code escaped the facility. Stuxnet infected over 100,000 machines across countries like Indonesia, India, and the U.S., burning four unpatched Windows zero-days and exposing unrelated networks to collateral risk [1], [2], [10], [15]. On top of that, Virtue Ethics and engineering codes like SECEPP focus on character and professional honesty [5], [12]. The weapon relied on active deception by spoofing normal sensor readouts to human operators while secretly destroying equipment [10]. Releasing an autonomous worm with flawed containment showed a deficit in the virtues of prudence and temperance, directly conflicting with professional obligations to protect public safety and disclose known hazards (SECEPP Clauses 1.03 and 1.04) [12]. Once the worm was discovered and decompiled, it essentially published a blueprint for attacking critical civilian utilities worldwide [15], [19].

Using the IEEE/SECEPP guidelines as a reference, we can find some comparisons and reflection on the ethical standing of Stuxnet, and attempt to apply them to our ethical theories.

1.02 - Public Interest: It is debatable whether stuxnet satisfied this clause, as one must consider the difference in damage that could have been caused versus what actually occurred. The US and Israel were concerned about nuclear weapons, with the potential for mass harm, and in doing so affected the public good of Iranian citizens and other unrelated innocents in the name of protecting the world.

1.03 - Software Safety: The software was designed to strategically damage equipment, and spread beyond the initial target. Even without this device going beyond its initial intended scope, it absolutely caused harm to the public and property.

3.10 - Testing and Review: Given that Stuxnet travelled beyond its originally intended targets, and began affecting infrastructure beyond the borders of Iran, it is plausible that insufficient testing was undergone related to the possibility of unrestricted spread.

Stuxnet was, at its core, an attack on a sovereign nation, and impeded on the rights and liberty of that nation. However, its ethical standing is a disputable point despite this, due to the possibility of massive negative consequences for the public if the attack was not completed. Stuxnet prevented the possible escalation of a nuclear arsenal in Iran, and in doing so prevented a country that did not have access to nuclear warheads from acquiring them.

We mention in our ethical analysis a rule that if followed, we believed would be beneficial to the most people. That rule was this: “Nations should use targeted, non-lethal sabotage to prevent new powers from acquiring nuclear weapons, because every new nation with access to a nuclear arsenal increases the risk of catastrophic use”. Even in consideration of analysis not pertaining to rule utilitarianism, this rule is a reasonable method of examining the ethics of the action the United States and Israel took. More nuclear weapons are inherently dangerous, as the more countries that have access to them the more risk there is for one deciding to set off armageddon. Therefore, something HAD to be done one way or another. In our discussion of act utilitarianism, we consider those alternate outcomes.

In one direction, we looked at diplomatic, non violent strategies like economic sanctions. These kinds of sanctions are not only often ineffective at short term prevention, but can cause ordinary citizens undue harm due to economic struggles influencing their daily lives.

On the opposite end of that spectrum, a full, boots on the ground military invasion of iran to prevent nuclear development was an unacceptable possibility. The loss of life on both sides would be catastrophic, and while it might prevent the development of a nuclear weapon, alternate options must be considered.

This leads us to our middle ground, a targeted cyber attack on Iranian nuclear infrastructure, designed to penetrate their defense then self replicate until their capability to produce nuclear devices was ruined. This did, of course, result in complications as the weapon spread more then intended, affecting unrelated parties outside of the original intended scope. However, this would be a more short term affliction than those caused by heavy economic sanctions and tariffs, and thus remains as the best option even in hindsight.

In our comparison with IEEE Guidelines, we could clearly see that the device did indeed not follow all the requirements set forth by IEEE, but while being a point to consider in ethical analysis, it is not the primary driving force in our decision. In conclusion, the Stuxnet attack on Iran’s nuclear program was an ethically permissible act, with their being room for discussion on ethical obligation. It prevented a nation from developing nuclear capability, and thus reduced the future risk for annihilation of the human race. When taken in this context, the unfortunate consequences of Stuxnet's unintended spread are, while terrible and a lesson for the future, an acceptable outcome overall.


% \section*{References}

% [1] J. Bentham, An Introduction to the Principles of Morals and Legislation. Oxford, U.K.: Clarendon Press, 1907. (Original work published 1789.)
% [2] J. S. Mill, Utilitarianism, 2nd ed., G. Sher, Ed. Indianapolis, IN, USA: Hackett, 2001.
% [3] B. Hooker, Ideal Code, Real World: A Rule-Consequentialist Theory of Morality. Oxford, U.K.: Oxford Univ. Press, 2000.
% [4] J. O. Urmson, “The interpretation of the moral philosophy of J. S. Mill,” The Philosophical Quarterly, vol. 3, no. 10, pp. 33–39, 1953.
% [5] Aristotle, Nicomachean Ethics, 2nd ed., T. Irwin, Trans. Indianapolis, IN, USA: Hackett, 1999.
% [6] I. Kant, Groundwork of the Metaphysics of Morals, M. Gregor, Trans. Cambridge, U.K.: Cambridge Univ. Press, 1997.
% [7] W. D. Ross, The Right and the Good. Oxford, U.K.: Clarendon Press, 1930.
% [8] T. Hobbes, Leviathan, R. Tuck, Ed. Cambridge, U.K.: Cambridge Univ. Press, 1996.
% [9] J. Locke, Two Treatises of Government, P. Laslett, Ed. Cambridge, U.K.: Cambridge Univ. Press, 1988.
% [10] N. Falliere, L. O. Murchu, and E. Chien, “W32.Stuxnet dossier,” Symantec Security Response, ver. 1.4, Feb. 2011. [Online]. Available: https://nsarchive.gwu.edu/document/21440-document-44
% [11] International Committee of the Red Cross, The Geneva Conventions of 12 August 1949. Geneva, Switzerland: ICRC. Accessed: Nov. 9, 2026. [Online]. Available: https://www.icrc.org/en/publication/0173-geneva-conventions-august-12-1949
% [12] IEEE-CS/ACM Joint Task Force on Software Engineering Ethics and Professional Practices, “Software engineering code of ethics and professional practice,” IEEE Computer Society. Accessed: Nov. 9, 2026. [Online]. Available: https://www.computer.org/education/code-of-ethics
% [13] M. A. Kamiński, “Operation ‘Olympic Games.’ Cyber-sabotage as a tool of American intelligence aimed at counteracting the development of Iran’s nuclear programme,” Security and Defence Quarterly, vol. 29, no. 2, pp. 63–71, 2020, doi: 10.35467/sdq/121974.
% [14] Kaspersky, “Stuxnet explained: What it is, who created it and how it works,” Kaspersky Resource Center. Accessed: Nov. 9, 2026. [Online]. Available: https://www.kaspersky.com/resource-center/definitions/what-is-stuxnet
% [15] P. K. Kerr, J. Rollins, and C. A. Theohary, “The Stuxnet computer worm: Harbinger of an emerging warfare capability,” Congressional Research Service, Washington, DC, USA, Rep. R41524, Dec. 9, 2010. [Online]. Available: https://www.congress.gov/crs_external_products/R/PDF/R41524/R41524.3.pdf
% [16] D. Kushner, “The real story of Stuxnet,” IEEE Spectrum, Feb. 26, 2013. [Online]. Available: https://spectrum.ieee.org/the-real-story-of-stuxnet
% [17] “Operation Olympic Games,” Wikipedia. Accessed: Nov. 9, 2026. [Online]. Available: https://en.wikipedia.org/wiki/Operation_Olympic_Games
% [18] “Stuxnet,” Wikipedia. Accessed: Nov. 9, 2026. [Online]. Available: https://en.wikipedia.org/wiki/Stuxnet
% [19] U.S. House of Representatives, Committee on Homeland Security, Subcommittee on Cybersecurity and Infrastructure Protection, “Fully operational: Stuxnet 15 years later and the evolution of cyber threats to critical infrastructure,” Washington, DC, USA, Jul. 22, 2025. [Online]. Available: https://homeland.house.gov/wp-content/uploads/2025/07/2025-07-22-CIP-Testimony.pdf
% [20] A. Gibney, Dir., Zero Days. Magnolia Pictures, 2016. [Film].

\begin{thebibliography}{99}

\bibitem{ref01_bentham1907}
J.~Bentham, \emph{An Introduction to the Principles of Morals and Legislation}. Oxford, U.K.: Clarendon Press, 1907. (Original work published 1789.)

\bibitem{ref02_mill2001}
J.~S. Mill, \emph{Utilitarianism}, 2nd ed., G.~Sher, Ed. Indianapolis, IN, USA: Hackett, 2001.

\bibitem{ref03_hooker2000}
B.~Hooker, \emph{Ideal Code, Real World: A Rule-Consequentialist Theory of Morality}. Oxford, U.K.: Oxford Univ. Press, 2000.

\bibitem{ref04_urmson1953}
J.~O. Urmson, ``The interpretation of the moral philosophy of J. S. Mill,'' \emph{The Philosophical Quarterly}, vol.~3, no.~10, pp.~33--39, 1953.

\bibitem{ref05_aristotle1999}
Aristotle, \emph{Nicomachean Ethics}, 2nd ed., T.~Irwin, Trans. Indianapolis, IN, USA: Hackett, 1999.

\bibitem{ref06_kant1997}
I.~Kant, \emph{Groundwork of the Metaphysics of Morals}, M.~Gregor, Trans. Cambridge, U.K.: Cambridge Univ. Press, 1997.

\bibitem{ref07_ross1930}
W.~D. Ross, \emph{The Right and the Good}. Oxford, U.K.: Clarendon Press, 1930.

\bibitem{ref08_hobbes1996}
T.~Hobbes, \emph{Leviathan}, R.~Tuck, Ed. Cambridge, U.K.: Cambridge Univ. Press, 1996.

\bibitem{ref09_locke1988}
J.~Locke, \emph{Two Treatises of Government}, P.~Laslett, Ed. Cambridge, U.K.: Cambridge Univ. Press, 1988.

\bibitem{ref10_falliere2011}
N.~Falliere, L.~O. Murchu, and E.~Chien, ``W32.Stuxnet dossier,'' Symantec Security Response, ver.~1.4, Feb. 2011. [Online]. Available: \url{https://nsarchive.gwu.edu/document/21440-document-44}

\bibitem{ref11_icrc1949}
International Committee of the Red Cross, \emph{The Geneva Conventions of 12 August 1949}. Geneva, Switzerland: ICRC. Accessed: Nov. 9, 2026. [Online]. Available: \url{https://www.icrc.org/en/publication/0173-geneva-conventions-august-12-1949}

\bibitem{ref12_ieeecs2026}
IEEE-CS/ACM Joint Task Force on Software Engineering Ethics and Professional Practices, ``Software engineering code of ethics and professional practice,'' IEEE Computer Society. Accessed: Nov. 9, 2026. [Online]. Available: \url{https://www.computer.org/education/code-of-ethics}

\bibitem{ref13_kaminski2020}
M.~A. Kamiński, ``Operation `Olympic Games.' Cyber-sabotage as a tool of American intelligence aimed at counteracting the development of Iran's nuclear programme,'' \emph{Security and Defence Quarterly}, vol.~29, no.~2, pp.~63--71, 2020, doi: 10.35467/sdq/121974.

\bibitem{ref14_kaspersky2026}
Kaspersky, ``Stuxnet explained: What it is, who created it and how it works,'' Kaspersky Resource Center. Accessed: Nov. 9, 2026. [Online]. Available: \url{https://www.kaspersky.com/resource-center/definitions/what-is-stuxnet}

\bibitem{ref15_kerr2010}
P.~K. Kerr, J.~Rollins, and C.~A. Theohary, ``The Stuxnet computer worm: Harbinger of an emerging warfare capability,'' Congressional Research Service, Washington, DC, USA, Rep. R41524, Dec. 9, 2010. [Online]. Available: \url{https://www.congress.gov/crs_external_products/R/PDF/R41524/R41524.3.pdf}

\bibitem{ref16_kushner2013}
D.~Kushner, ``The real story of Stuxnet,'' \emph{IEEE Spectrum}, Feb. 26, 2013. [Online]. Available: \url{https://spectrum.ieee.org/the-real-story-of-stuxnet}

\bibitem{ref17_wiki_olympic2026}
``Operation Olympic Games,'' Wikipedia. Accessed: Nov. 9, 2026. [Online]. Available: \url{https://en.wikipedia.org/wiki/Operation_Olympic_Games}

\bibitem{ref18_wiki_stuxnet2026}
``Stuxnet,'' Wikipedia. Accessed: Nov. 9, 2026. [Online]. Available: \url{https://en.wikipedia.org/wiki/Stuxnet}

\bibitem{ref19_ushouse2025}
U.S. House of Representatives, Committee on Homeland Security, Subcommittee on Cybersecurity and Infrastructure Protection, ``Fully operational: Stuxnet 15 years later and the evolution of cyber threats to critical infrastructure,'' Washington, DC, USA, Jul. 22, 2025. [Online]. Available: \url{https://homeland.house.gov/wp-content/uploads/2025/07/2025-07-22-CIP-Testimony.pdf}

\bibitem{ref20_gibney2016}
A.~Gibney, Dir., \emph{Zero Days}. Magnolia Pictures, 2016. [Film].

\end{thebibliography}




% \bibliographystyle{IEEEtran}
% \bibliography{refs}
